% 2015 年重庆卷（文史类）第 5 题：某几何体的三视图（依固定扫描页重绘，只画拓扑与标号，不按比例）
\begin{tikzpicture}[>=stealth,scale=1.0]
  % 正视图：边长 2 的正方形，左下接直角三角形（水平边 1、竖直边 1）
  \draw (0,0) rectangle (2,2);
  \draw (-1,1)--(0,1)--(0,0)--cycle;
  % 正视图尺寸标注
  \draw[<->] (-1,1.45)--(0,1.45);
  \node at (-0.5,1.72) {$1$};
  \draw[<->] (0.28,1)--(0.28,2);
  \node at (0.62,1.5) {$1$};
  \draw[<->] (0.28,0)--(0.28,1);
  \node at (0.62,0.5) {$1$};
  \node at (0.5,-0.55) {正视图};
  % 左视图：半径 1 的圆，水平直径
  \draw (4.5,1) circle (1);
  \draw (3.5,1)--(5.5,1);
  \node at (4.5,-0.55) {左视图};
  % 俯视图：边长 2 的正方形，左侧接等腰三角形（底 2、高 1）
  \draw (0,-4) rectangle (2,-2);
  \draw (-1,-3)--(0,-2)--(0,-4)--cycle;
  % 俯视图尺寸标注
  \draw[<->] (2.35,-2)--(2.35,-4);
  \node at (2.72,-3) {$2$};
  \draw[<->] (0,-4.4)--(2,-4.4);
  \node at (1,-4.7) {$2$};
  \node at (0.5,-5.25) {俯视图};
\end{tikzpicture}
